\subsection{A genuine generic graph retains the SAME original function field}

\begin{definition}[Actual coefficient homomorphism of a dominant field-valued point]
\label{scheme:generic-field-point-hom}
\lean{EnvelopeProperty.SchemeGenericFieldGraph.fieldHom}
\leanok
For an integral scheme $Y$ and an ACTUAL dominant map
$g:\operatorname{Spec}K\to Y$, define the actual field homomorphism
$k(Y)\to K$ by the native stalk map at the actual image of the closed point,
transporting only along its proved equality with the genuine generic point.
\end{definition}
\begin{proof}
The field spectrum has one point; its original image is generic by true
dominance. Thus the native closed-point stalk map has, after actual native
stalk transport, exactly the actual generic stalk as source.
\end{proof}

\begin{theorem}[The entire original field point is retained]
\label{scheme:generic-field-point-hom-diagram}
\lean{EnvelopeProperty.SchemeGenericFieldGraph.fieldHom_fromSpecStalk}
\leanok
\uses{scheme:generic-field-point-hom}
The spectrum of this SAME actual coefficient homomorphism, followed by the
native generic-stalk morphism, equals the SAME original $g$ as an entire
scheme morphism.
\end{theorem}
\begin{proof}
The genuine native stalk-transport triangle moves the generic stalk to the
actual image stalk. The original closed-point-to-stalk factorization then
recovers the whole original map $g$, including its actual section maps.
\end{proof}

\begin{theorem}[An actual generic graph has exactly the original function field]
\label{scheme:generic-graph-original-function-field}
\lean{EnvelopeProperty.SchemeGenericFieldGraph.fieldHom_bijective_of_genericProjection,EnvelopeProperty.SchemeGenericFieldGraph.functionFieldEquiv}
\leanok
\uses{scheme:generic-field-point-hom-diagram}
Let $X,Y$ be actual integral schemes and let
$g:\operatorname{Spec}k(X)\to Y$ be an actual dominant point. If an actual
whole projection $p:Y\to X$ satisfies
$g\circ p=X.\mathrm{fromSpecStalk}(\eta_X)$, the SAME actual native
coefficient homomorphism $k(Y)\to k(X)$ is bijective and defines a genuine
field isomorphism. No birationality or function-field comparison is assumed.
\end{theorem}
\begin{proof}
The genuine generic point of $X$ lies in the actual image of $p$, so $p$ is
dominant. Applying the actual native stalk map of $p$ at $\eta_Y$, with
its proved generic-point equality, gives an actual map
$\operatorname{Spec}k(Y)\to X$ equal to the native generic morphism of $Y$
followed by $p$. Regard this as a dominant $k(Y)$-valued point of $X$ and
form its genuine coefficient homomorphism $k(X)\to k(Y)$.

Compose the two proven whole generic diagrams and use the original equation
$g\circ p=X.\mathrm{fromSpecStalk}(\eta_X)$. Cancelling the native
monomorphism from the generic stalk of $X$ proves that the two spectrum
maps compose to the identity of $\operatorname{Spec}k(X)$. Full faithfulness
of actual $\operatorname{Spec}$ then proves that the original coefficient
homomorphism $k(Y)\to k(X)$ has a right inverse. It is injective because
it is a field homomorphism, and surjective by this right inverse, giving the
claimed SAME-field isomorphism.
\end{proof}
